\chapter{Background and Related Work}

\section{Federated Learning}
\ \ \ \ \, Federated learning considers the optimization of a global model over data
distributed across multiple clients. If client $n$ contributes to local data
$\mathcal{D}_n$, the global objective can be written as
\begin{equation}
    \min_w F(w) =
    \sum_{n=1}^{N} p_n F_n(w),
    \quad
    F_n(w) =
    \frac{1}{|\mathcal{D}_n|}
    \sum_{(x,y)\in\mathcal{D}_n}
    \ell(w;x,y),
\end{equation}
where $p_n$ is often proportional to the local data size $|\mathcal{D}_n|$. FedAvg trains this
objective by alternating between local client optimization and server-side
averaging \citep{mcmahan2017communication}. The method is simple and
communication-efficient comparing to fully synchronous distributed SGD, but
it still requires explicit client uploads in each round.

\section{Analog Over-the-Air Aggregation}
\ \ \ \ \, OTA-FL replaces separate digital uploads with analog aggregation in the
wireless channel. In an idealized channel, simultaneous client transmissions
sum at the receiver, so the server directly observes an aggregate update, limiting
the required multiple-access resource's sensitivity to
the number of clients. However, in reality the observed aggregate is usually noisy. A simplified
channel model is
\begin{equation}
    y^t = \sum_{n=1}^{N} \Delta_n^t + \xi^t,
\end{equation}
As equation $(1.2)$ shows, the server uses $\tilde{g}_t = y_t/N$ as the update direction. In the
AWGN case where $\alpha=2$, $\xi^t \sim \mathcal{N}(0,\sigma^2 I)$. In the heavier
channel model used by this project, $\xi^t$ follows a symmetric
$\alpha$-stable distribution.

The original adaptive OTA-FL paper includes more general channel effects,
e.g. fading coefficients and symmetric $\alpha$-stable interference
\citep{wang2024adaptive}. This project, however, does not implement fading or explicit
power-control mechanisms, but the $\alpha$-stable interference
axis. This is the key step beyond an AWGN-only simulator because
$\alpha$-stable noise captures rare but large impulsive channel perturbations.

\section{Adaptive Gradient Methods}

\ \ \ \, Adaptive optimizers adapts the learning rate separately for each varaibles by 
the historical gradients. Here we introduce AdaGrad and Adam.

AdaGrad accumulates squared gradients and scales each coordinate's
learning rate inversely proportional to the square root of its accumulated
magnitude \citep{duchi2011adaptive}. Adam combines first-moment smoothing,
second-moment exponential averaging and bias correction
\citep{kingma2015adam}. In centralized learning, these mechanisms are usually
motivated by sparse gradients, noisy mini-batches and heterogeneous
coordinate scales.

In OTA-FL, the same mechanisms have an additional interpretation. If channel
noise makes the aggregate update unusually large in some coordinate, the
second-moment state for that coordinate increases. The denominator of the
adaptive update then becomes larger, reducing the effective step size. This
creates an implicit feedback loop from observed noisy update magnitude to
server learning rate. The project tests whether this feedback loop is visible
in training curves and ablation studies.

\section{Positioning Relative to the Reference Paper}
\ \ \ \ \, The reference paper proposes adaptive federated learning over the air and
studies AdaGrad-like and Adam-like server updates with theoretical analysis
under $\alpha$-stable interference \citep{wang2024adaptive}. It evaluates
multiple image, character classification tasks and includes studies over
tail index, client count and data heterogeneity.

The present project uses that paper as the conceptual source but narrows the
deliverable. The implemented framework supports the core OTA-FL training loop,
$\alpha$-stable noise based on the Chambers-Mallows-Stuck method
\citep{chambers1976method}, MAC robust pre-processing and adaptive server
optimizers. The main evidence is the alpha tail-index ablation, with AWGN
reported as the $\alpha=2$ endpoint rather than the whole scope.
